\documentclass[../../main.tex]{subfiles}
\begin{document}


{\bf [解]}
簡単のため$c=\cos x, s=\sin x$とおく．

\paragraph{(1)}
新しく
\begin{align}
 F(x)=p(x)c-q(x)s \label{eq:1}
\end{align}
とおくと，全ての$x$について$F(x)=0$である場合を考えれば良い．
$p(x), q(x)$ が恒等的に$0$でないと仮定する．$p(x), q(x)$は$3$次以下なので高々$3$個の$x$でしか$p(x)=q(x)=0$とならない．
それ以外の$x$の時，\cref{eq:1}に対して三角関数の合成を行って
\begin{align}
   F(x)=\sqrt{p^2(x)+q^2(x)}\sin(x+\alpha) \quad \left(\alpha\text{は}\ \tan\alpha=-\dfrac{P(x)}{Q(x)}\ \text{をみたす}\ -\pi/2<\alpha<\pi/2\text{の数}\right) 
\end{align}
と書ける．しかしこの右辺が任意の $x$ で0になることはないので $F(x)=0$に矛盾する．
よって背理法により $p(x), q(x)$ は恒等的に$0$ である．

\paragraph{(2)}
与えられた方程式
\begin{align}
  P(x)\cdot c+\int_0^xQ(t)\sin t\,dt=(x^2+2x+3)s \quad \cdots \text{①} \label{eq:2}
\end{align}
を考える．まず$x=0$を\cref{eq:2}に代入すると
\begin{align}
 P(0) = 0 \label{eq:3}
\end{align}
である．次に\cref{eq:2}の両辺を $x$ で微分して
\begin{align}
\begin{split}
  -P(x)s+P'(x)c+Q(x)\cdot s&=(2x+2)s+(x^2+2x+3)c \\
  \therefore 
  [Q(x)-P(x)-(2x+2)]s&=(x^2+2x+3-P'(x))c
\end{split}
\end{align}
を得る．ここで$P(x), Q(x)$が$3$次以下なので$Q(x)-P(x)-(2x+2)$，$x^2+2x+3-P'(x)$も$3$次以下である．
よって(1)の結果が使えて
\begin{align}
  \begin{cases}\displaystyle
  Q(x)-P(x)=2x+2 \\
  P'(x)=x^2+2x+3
  \end{cases} \label{eq:4} 
\end{align}
である．
\cref{eq:4}の第二式の両辺を積分して\cref{eq:3}と合わせて
\begin{align}
   P(x)=\dfrac{1}{3}x^3+x^2+3x  \label{eq:5}
\end{align}
となる．\cref{eq:4}の第一式にこれを代入して
\begin{align}
  Q(x)=\dfrac{1}{3}x^3+x^2+5x+2 \label{eq:6}
\end{align}
となる．\cref{eq:5,eq:6}はいずれも$3$次以下の多項式でかつ\cref{eq:2}を満たしており，よって求める関数は
\begin{align}
 \begin{cases}
  P(x)=\dfrac{1}{3}x^3+x^2+3x \\
  Q(x)=\dfrac{1}{3}x^3+x^2+5x+2
 \end{cases}
\end{align}
である．

\newpage

\end{document}
